% Ad Familiares 2.16 (ad-familiares-02-16)
% Source: https://raw.githubusercontent.com/PerseusDL/canonical-latinLit/master/data/phi0474/phi056/phi0474.phi056.perseus-lat1.xml
% Fetched by scripts/fetch_latin.py

% Dateline (Perseus): Scr. in Cumano iv Non. Mai. a. 705 (49).

\ciceroLetterOpener{M. CICERO IMP. S. D. M. CAELIO.}

\ciceroSection{1} Magno dolore me adfecissent tuae litterae, nisi iam et ratio ipsa depulisset omnis molestias et diuturna desperatione rerum obduruisset animus ad dolorem novum. sed tamen, qua re acciderit ut ex meis superioribus litteris id suspicarere, quod scribis, nescio; quid enim in illis fuit praeter querelam temporum? quae non meum animum magis sollicitum habent quam tuum. nam non eam cognovi aciem ingeni tui, quod ipse videam, te id ut non putem videre; illud miror, adduci potuisse te, qui me penitus nosse deberes, ut existimares aut me tam improvidum, qui ab excitata fortuna ad inclinatam et prope iacentem desciscerem, aut tam inconstantem, ut conlectam gratiam florentissimi hominis effunderem a meque ipse deficerem et, quod initio semperque fugi, civili bello interessem. quod est igitur meum 'triste consilium'?

\ciceroSection{2} ut discederem fortasse in aliquas solitudines. Nosti enim non modo stomachi mei, cuius tu similem quondam habebas, sed etiam oculorum in hominum insolentium indignitate fastidium. accedit etiam molesta haec pompa lictorum meorum nomenque imperi, quo appellor. eo si onere carerem, quamvis parvis Italiae latebris contentus essem; sed incurrit haec nostra laurus non solum in oculos, sed iam etiam in voculas malevolorum. quod cum ita esset, nil tamen umquam de profectione nisi vobis approbantibus cogitavi. sed mea praediola tibi nota sunt; in his mihi necesse est esse, ne amicis molestus sim. quod autem in maritimis facillime sum, moveo non nullis suspicionem velle me navigare; quod tamen fortasse non nollem, si possem ad otium; nam ad bellum quidem qui convenit? praesertim contra eum, cui spero me satis fecisse, ab eo, cui iam satis fieri nullo modo potest.

\ciceroSection{3} deinde sententiam meam tu facillime perspicere potuisti iam ab illo tempore, cum in Cumanum mihi obviam venisti. non enim te celavi sermonem T. Ampi; vidisti, quam abhorrerem ab urbe relinquenda, cum audissem; nonne tibi adfirmavi quidvis me potius perpessurum quam ex Italia ad bellum civile exiturum? quid ergo accidit, cur consilium mutarem? nonne omnia potius, ut in sententia permanerem? credas hoc mihi velim, quod puto te existimare, me ex his miseriis nihil aliud quaerere nisi ut homines aliquando intellegant me nihil maluisse quam pacem, ea desperata nihil tam fugisse quam arma civilia. huius me constantiae puto fore ut numquam paeniteat. etenim memini in hoc genere gloriari solitum esse familiarem nostrum, Q Hortensium, quod numquam bello civili interfuisset; hoc nostra laus erit inlustrior, quod illi tribuebatur ignaviae, de nobis id existimari posse non arbitror.

\ciceroSection{4} nec me ista terrent, quae mihi a te ad timorem fidissime atque amantissime proponuntur. nulla est enim acerbitas, quae non omnibus hac orbis terrarum perturbatione impendere videatur; quam quidem ego a re publica meis privatis et domesticis incommodis libentissime vel istis ipsis, quae tu me mones ut caveam, redemissem.

\ciceroSection{5} filio meo, quem tibi carum esse gaudeo, si erit ulla res publica, satis amplum patrimonium relinquam in memoria nominis mei; sin autem nulla erit, nihil accidet ei separatim a reliquis civibus. nam quod rogas ut respiciam generum meum, adulescentem optimum mihique carissimum, an dubitas, qui scias quanti cum illum tum vero Tulliam meam faciam, quin ea me cura vehementissime sollicitet, et eo magis, quod in communibus miseriis hac tamen oblectabar specula, Dolabellam meum vel potius nostrum fore ab iis molestiis, quas liberalitate sua contraxerat, liberum? velim quaeras, quos ille dies sustinuerit, in urbe dum fuit, quam acerbos sibi, quam mihimet ipse socero non honestos.

\ciceroSection{6} itaque neque ego hunc Hispaniensem casum exspecto, de quo mihi exploratum est ita esse, ut tu scribis, neque quicquam astute cogito. si quando erit civitas, erit profecto nobis locus; sin autem non erit, in easdem solitudines tu ipse, ut arbitror, venies, in quibus nos consedisse audies. sed ego fortasse vaticinor et haec omnia meliores habebunt exitus. recordor enim desperationes eorum, qui senes erant adulescente me. Eos ego fortasse nunc imitor et utor aetatis vitio. velim ita sit; sed tamen—.

\ciceroSection{7} togam praetextam texi Oppio puto te audisse; nam Curtius noster dibaphum cogitat, sed eum infector moratur. hoc aspersi, ut scires me tamen in stomacho solere ridere. de re Dolabellae quod scripsi, suadeo videas, tamquam si tua res agatur. extremum illud erit: nos nihil turbulenter, nihil temere faciemus; te tamen oramus, quibuscumque erimus in terris, ut nos liberosque nostros ita tueare, ut amicitia nostra et tua fides postulabit.
